\begin{textsample}{POS Dim 4 – vem\_ed\_31 – Score 1.00 – vem\_ed\_31/t160.txt}  \label{ex:f4_pos_147}
Aos Leitores Osvaldo Succi Jr. - Coordenador dos PCIs A equipe de Apoio à Internacionalização do Ensino Superior da Divisão de Extensão e Pesquisa no Ensino Superior da Coordenadoria Geral de Ensino Superior de Graduação do CPS marca presença na movimentada agenda de eventos relacionados aos Intercâmbios Virtuais e aprendizagem de línguas estrangeiras.
No VII CBTecLE, realizado dias 11 e 12 de setembro de 2025 na Fatec Praia Grande, a equipe ofereceu um workshop sobre as ferramentas de Inteligência Artificial na gestão de Projetos Colaborativos Internacionais (PCIs).
Além disso, cinco PCIs foram apresentados no congresso por professores das Fatecs Barueri, Bragança Paulista, Ipiranga, São Paulo e Taquaritinga.
No Colóquio de Educação Internacional para o Sul Global, evento \textbf{online} realizado dia 2 de setembro, Neusa Haruka Sezaki Gritti e Patrícia Patrício compartilharam suas experiências com PCIs realizados com China e Portugal, respectivamente.
Regiane Camargo apresentou trabalhos sobre os projetos COIL (Collaborative Online International Learning) do CPS en español em dois eventos relacionados à aprendizagem de língua espanhola: I Congreso Internacional sobre Productividad Investigativa en el Aprendizaje de Lenguas, realizado \textbf{online} em 17 de setembro, e o Encuentro de Español para Fines Específicos, que ocorreu no Instituto Cervantes de São Paulo em 27 de setembro.
No SUNY COIL Community Webinar, em 17 de setembro, debati com Luciane Stallivieri (UFSC) como os projetos COIL podem preparar profissionais para empregos que ainda não existem.
A importância dos PCIs para a qualificação profissional contemporânea também é discutida no Artigo de Opinião assinado pelas professoras Kátia Cristina Galatti, Talita Botelho Nunes e Vanessa Amaro Vieira, da Fatec Taquaritinga.
Boa leitura!

% matched lemmas: onlinar
\end{textsample}
